\documentclass{beamer}

\usepackage{hyperref}
\usepackage[utf8]{inputenc}
\usepackage[T1]{fontenc}
\usepackage[spanish]{babel}

% --- Paquetes para diseño ---
\usepackage{latexsym,amsmath,xcolor,multicol,booktabs,calligra}
\usepackage{graphicx,listings,stackengine,tikz}

% --- Definición de Cajas Personalizadas ---
\newenvironment{cajaazul}[1]{%
    \setbeamercolor{block title}{bg=cyan!70!black, fg=white}
    \setbeamercolor{block body}{bg=cyan!10!white, fg=black}
    \begin{block}{#1}}{\end{block}}

\newenvironment{cajaverde}[1]{%
    \setbeamercolor{block title}{bg=teal!80!black, fg=white}
    \setbeamercolor{block body}{bg=teal!10!white, fg=black}
    \begin{block}{#1}}{\end{block}}

\newenvironment{cajaranja}[1]{%
    \setbeamercolor{block title}{bg=orange!90!black, fg=white}
    \setbeamercolor{block body}{bg=orange!10!white, fg=black}
    \begin{block}{#1}}{\end{block}}

% --- DATOS DEL PROYECTO ---
\author{Juan Pedro García Sanz \and Pablo Revuelto de Miguel \and \\ Adolfo Peña Marín}
\title{Práctica 1 --- Simulación del Agilex Tracer}
\subtitle{Gazebo Harmonic + ROS\,2 Jazzy}
\institute{Robótica Inteligente \\ Universidad Loyola Andalucía}

% El logo se añade aquí, justo encima de la fecha con un salto de línea
\date{\includegraphics[height=1cm]{logo.png} \\[0.2cm] \today}

% --- TU TEMA PERSONALIZADO ---
\usepackage{CNU}
\usefonttheme{professionalfonts}
\usepackage{lmodern}

% --- Colores para código ---
\definecolor{deepblue}{rgb}{0,0,0.5}
\definecolor{deepred}{rgb}{0.6,0,0}
\definecolor{deepgreen}{rgb}{0,0.5,0}
\definecolor{halfgray}{gray}{0.55}

\newcommand{\completar}[1]{\textcolor{halfgray}{\textit{[#1]}}}

\begin{document}

% --- PORTADA ---
\begin{frame}
    \titlepage
\end{frame}

% ============================================================
% ÍNDICE
% ============================================================
\begin{frame}{Índice}
    \tableofcontents[hideallsubsections]
\end{frame}

% ============================================================
% Objetivo y alcance
% ============================================================
\section{Objetivo y Alcance}
\begin{frame}{Qué hemos construido}

\begin{cajaazul}{Entregables de la práctica}
\begin{itemize}\setlength{\itemsep}{0pt}
    \item \textbf{Mundo SDF}: física, fricción regulable y obstáculos.
    \item \textbf{Robot}: Agilex Tracer (diferencial) con IMU, LiDAR 2D y cámara.
    \item \textbf{Integración ROS 2}: puente de teleop, odometría, TF y RViz 2.
    \item \textbf{Experimento D1}: análisis de paso de solver y fricción.
\end{itemize}
\end{cajaazul}

\begin{cajaverde}{Entorno de software}
\textbf{ROS 2 Jazzy} + \textbf{Gazebo Harmonic}. Plugins migrados de \texttt{ignition-gazebo-*} a \texttt{gz-sim-*}.
\end{cajaverde}

\end{frame}

% ============================================================
% Parte A: Mundo y Robot
% ============================================================
\section{Mundo SDF y Robot (URDF)}
\begin{frame}{Configuración del Mundo SDF}

\begin{cajaazul}{Plugins y Física (gz-sim)}
\begin{itemize}\setlength{\itemsep}{0pt}
    \item \textbf{Sistemas}: \texttt{physics}, \texttt{user-commands}, \texttt{scene-broadcaster}, \texttt{sensors} (render Ogre2), \texttt{imu} y \texttt{contact}.
    \item \textbf{Física}: bloque \texttt{<physics>} regulable (base del experimento D1).
    \item \textbf{Suelo}: fricción configurable vía \texttt{<surface><friction><ode>}.
\end{itemize}
\end{cajaazul}

\begin{cajaverde}{Obstáculos, cámara inicial y validación}
Paredes (4 ext. + 3 int.), cajas, mesa y paso estrecho (estáticos con colisión). Bloque \texttt{<gui><camera>} con pose inicial elevada/diagonal. Validado mediante comando \texttt{gz sdf -k} $\rightarrow$ \textit{Valid.}
\end{cajaverde}

\end{frame}

% ============================================================
% Parte B: Robot
% ============================================================
\begin{frame}{Estructura del Robot}

\begin{cajaazul}{Distribución física (6 ruedas y chasis)}
\begin{itemize}\setlength{\itemsep}{0pt}
    \item \textbf{Actuación}: 2 ruedas motrices centrales (\texttt{diff-drive}).
    \item \textbf{Estabilidad}: 4 castors pasivos en las esquinas como anti-vuelco.
    \item \textbf{Estructura}: chasis (\texttt{base\_link}) y soporte para sensores.
\end{itemize}
\end{cajaazul}

\begin{cajaverde}{Árbol de transformadas (TF)}
\textbf{Estáticas}: publicadas por \texttt{robot\_state\_publisher}.\\
\textbf{Dinámicas}: publicadas por el plugin de juntas del robot para RViz 2.
\end{cajaverde}

\end{frame}

\begin{frame}{Correcciones en el URDF Original}

\begin{cajaranja}{Bug crítico: eje invertido en rueda izquierda}
Llevaba \texttt{rpy="3.14 0 0"} $\rightarrow$ invertía su eje físico de rotación. Al avanzar, las ruedas giraban en sentido opuesto y el robot rotaba sobre sí mismo.
\end{cajaranja}

\begin{cajaverde}{Correcciones geométricas y dinámicas}
\begin{itemize}\setlength{\itemsep}{0pt}
    \item \textbf{Medidas reales} del Tracer: separación ($0.34$\,m) y radio ($0.059$\,m). Nota: chasis 30\,kg y diámetro 0.118\,m se apartan del rango sugerido como ``p.\,ej.'' en B1 (5--15\,kg, 0.2\,m), pero corresponden al robot real.
    \item \textbf{Contactos}: rodadura con cilindro primitivo. Castor con $\mu=0.05$ para permitir el deslizamiento lateral en los giros.
    \item \textbf{Estabilidad}: altura de ruedas calibrada ($z=-0.085$\,m) e inercias coherentes (caja para chasis, cilindro para ruedas).
\end{itemize}
\end{cajaverde}

\end{frame}

% ============================================================
% Parte C: Sensores e Integración
% ============================================================
\section{Sensores e Integración ROS 2}
\begin{frame}{Sensores Embarcados}

\begin{cajaazul}{Especificaciones técnicas}
\centering\small
\begin{tabular}{@{}lll@{}}
\toprule
\textbf{Sensor} & \textbf{Frame} & \textbf{Tópico ROS} \\
\midrule
IMU (100\,Hz)                       & \texttt{imu\_link}    & \texttt{/imu} \\
LiDAR 2D (FOV $270^\circ$, 12\,m)  & \texttt{lidar\_link}  & \texttt{/scan} \\
Cámara ($640\times480$, 30\,Hz)     & \texttt{camera\_link} & \texttt{/camera/image\_raw} \\
Contacto (evento al golpear)        & \texttt{base\_link}   & \texttt{/contacts} \\
\bottomrule
\end{tabular}
\end{cajaazul}

\begin{cajaverde}{Verificación y nota extra}
\texttt{frame\_id} correctos en todos los sensores. LiDAR y cámara estables en RViz 2. El sensor de \textbf{contacto} cubre el objetivo 3 del enunciado y la extensión \#2 (sección 14).
\end{cajaverde}

\end{frame}

\begin{frame}{Puente e Integración ROS 2}

\begin{cajaazul}{Comunicación bidireccional (\texttt{ros\_gz\_bridge})}
\begin{itemize}\setlength{\itemsep}{0pt}
    \item \textbf{Actuación}: puente de velocidad en \texttt{/cmd\_vel}.
    \item \textbf{Feedback}: odometría (\texttt{/odom}), estados de juntas (\texttt{/joint\_states}), TF, IMU (\texttt{/imu}), LiDAR (\texttt{/scan}), cámara y \texttt{/contacts}.
    \item \textbf{Control}: teleoperación por teclado y visualización completa en RViz 2.
\end{itemize}
\end{cajaazul}

\begin{cajaverde}{Corrección en el Lanzador (ROS 2 Jazzy)}
Envoltura de la descripción en \texttt{ParameterValue} para corregir el fallo de lectura YAML de Jazzy. Exportación automatizada de \texttt{GZ\_SIM\_RESOURCE\_PATH}.
\end{cajaverde}

\end{frame}

% ============================================================
% Parte D: Experimento físico (D1)
% ============================================================
\section{Experimento Físico (D1)}
\begin{frame}{Experimento D1: Planificación}

\begin{cajaazul}{Configuraciones de la simulación}
\centering\small
\begin{tabular}{@{}llcc@{}}
\toprule
\textbf{Corrida} & \textbf{Hipótesis} & $\mu$ Suelo & Paso (s) \\
\midrule
E1 fricción baja & Patina en giros & 0.05 & 0.001 \\
E2 fricción alta & Giro nominal estable & 1.0  & 0.001 \\
E3 paso grande & Inestabilidad del solver & 1.0  & 0.01  \\
\bottomrule
\end{tabular}
\end{cajaazul}

\begin{cajaverde}{Trayectoria de control idéntica}
Avance recto (3\,s) seguido de un giro a la izquierda (5\,s).
\end{cajaverde}

\end{frame}

\begin{frame}{Experimento D1: Resultados}

\begin{cajaazul}{Tabla comparativa}
\centering\footnotesize
\begin{tabular}{@{}lccc@{}}
\toprule
\textbf{Corrida} & \textbf{Patina} & \textbf{Estable} & \textbf{Sensores} \\
\midrule
E1 --- fricción baja ($\mu{=}0.05$) & No & Sí & OK \\
E2 --- fricción alta ($\mu{=}1.0$)  & No & Sí & OK \\
E3 --- paso grande (step$=0.01$)    & No & Sí & OK \\
\bottomrule
\end{tabular}
\end{cajaazul}

\begin{cajaranja}{Conclusiones físicas}
\begin{itemize}\setlength{\itemsep}{0pt}
    \item \textbf{E1 vs E2}: E2 describe un arco más cerrado — con mayor agarre el giro es más preciso.
    \item \textbf{E3}: paso grande no inestabiliza porque la dinámica del robot es suave.
\end{itemize}
\end{cajaranja}

\end{frame}

% ============================================================
% Anexo: migración Classic -> Sim
% ============================================================
\section{Anexo y Conclusiones}
\begin{frame}{Anexo: Migración Classic a Sim (Harmonic)}

\begin{cajaazul}{Cambios en Plugins y Mallas}
\begin{itemize}\setlength{\itemsep}{0pt}
    \item \textbf{Sistemas}: migración de plugins antiguos a \texttt{gz-sim-*-system}. Control diff-drive directo en el bloque \texttt{<gazebo>}.
    \item \textbf{Visuales}: mallas resueltas vía \texttt{model://} con \texttt{GZ\_SIM\_RESOURCE\_PATH} apuntando al directorio raíz de modelos.
\end{itemize}
\end{cajaazul}

\begin{cajaverde}{Tópicos y Lanzadores}
\begin{itemize}\setlength{\itemsep}{0pt}
    \item \textbf{Puente}: mapeo explícito de tópicos a través de \texttt{ros\_gz\_bridge}.
    \item \textbf{Lanzador en Jazzy}: uso de \texttt{ParameterValue} para la carga de descripción y corrección de errores de análisis YAML.
\end{itemize}
\end{cajaverde}

\end{frame}

\begin{frame}{Conclusiones}

\begin{cajaazul}{Logros clave de la práctica}
\begin{itemize}\setlength{\itemsep}{0pt}
    \item \textbf{Entorno robusto}: simulación totalmente funcional en ROS 2 Jazzy y Gazebo Harmonic.
    \item \textbf{Correcciones URDF}: resueltos problemas de inversión de eje, dimensiones y fricción de las ruedas.
    \item \textbf{Experimento físico}: comprobación del impacto de la fricción y el paso del solver.
    \item \textbf{Migración documentada}: guía práctica de Classic a Sim.
\end{itemize}
\end{cajaazul}

\begin{center}
\Large \textbf{¿Preguntas?}
\end{center}

\end{frame}

\end{document}
